% ===========================================================================
%  sviluppatore/07-marketplace/04-artefatti.tex
% ===========================================================================
\section{Pubblicare e installare}\label{sez:pubblicare}

\subsection{Nessuna moderazione, un filtro all'ingresso}

\begin{codiceInterno}{python}{La scelta, scritta dove si applica}{dashboard/community.py:1537}
        # CHI PUO' PUBBLICARE
        # Solo amministratori e abbonati, anche per gli artefatti
        # gratuiti. Il motivo non e' il denaro: e' che chi pubblica
        # mette il proprio nome accanto a una strategia che altri
        # faranno girare sui propri soldi. Un filtro all'ingresso vale
        # piu' di qualunque moderazione a valle.
        #
        # Chi non e' abbonato resta un utente pieno: sfoglia, scarica,
        # compra e usa. Non puo' soltanto mettere roba in vetrina.
\end{codiceInterno}

\noindent
Un artefatto pubblicato va online subito: nessuna coda di approvazione. La
difesa è a monte --- serve un abbonamento --- e a valle le contestazioni
(\vediSez{pagamenti}), non in mezzo.

\subsection{I conflitti si dicono prima}

\rotta{GET}{/api/marketplace/<id>/conflicts} risponde a «che cosa esiste già con
lo stesso nome». Si chiama \emph{prima} di installare, e serve a non scoprire
dopo che un'installazione ha scavalcato una strategia su cui si stava
lavorando. Il nome canonico di un artefatto è il nome della sua cartella
(\vediSez{due-forme}): due artefatti con lo stesso nome sono la stessa
cartella.

\subsection{La pubblicazione automatica}

\file{dashboard/autopubblica.py} gira in un thread dall'avvio e pubblica da
sé gli artefatti che superano una soglia.

\begin{codiceInterno}{python}{Che cosa fa}{dashboard/autopubblica.py}
COSA FA
Ogni volta che una strategy viene scritta, e ogni volta che un project viene
creato, si guarda lo Sharpe out-of-sample del walk-forward. Sopra 2.0 e non
ancora in vetrina, l'artefatto va online GRATIS da solo. Per un project la
regola e' piu' articolata (vedi `valuta_progetto`): dipende da quante strategie
lo compongono e da cosa e' gia' pubblicato.
\end{codiceInterno}

\subsubsection{Il file di tracciamento}

Senza, ogni avvio rifarebbe gli stessi controlli e --- peggio ---
riproverebbe le stesse pubblicazioni. Il file registra, per ogni artefatto,
se è stato visionato e con quale conclusione. È un JSON indentato, leggibile
a mano: non un formato interno.

\subsubsection{Quando una voce viene riesaminata}

Tre casi, e solo questi:

\begin{enumerate}
  \item non c'è (artefatto nuovo, o file di tracciamento appena nato);
  \item \textbf{la conclusione non era definitiva} --- mancava il backtest,
        mancava la rete, nessuno era connesso;
  \item l'impronta del codice o quella dei risultati è cambiata.
\end{enumerate}

\begin{perche}{La distinzione fra un ostacolo e un giudizio}
Il commento del modulo la enuncia meglio di qualunque parafrasi: gli
ostacoli «non sono giudizi sull'artefatto: sono ostacoli, e un ostacolo che
passa deve poter cambiare l'esito. Una soglia non superata, invece, è un
giudizio: resta tale finché il codice o i risultati non cambiano».

È il criterio per decidere se un esito negativo vada riprovato o meno, e si
applica ben oltre questo modulo: ogni volta che un sistema registra un
«no», la domanda è se quel no dipendeva dall'oggetto o dalle circostanze.
\end{perche}

\begin{attenzione}{Un artefatto può finire in vetrina da solo}
È un comportamento attivo che agisce sul mondo esterno senza che nessuno lo
chieda in quel momento. Chi lavora su \file{autopubblica.py} tenga presente
che un errore qui non produce un messaggio: produce una pubblicazione. Il
file di tracciamento, essendo leggibile a mano, è lo strumento per
verificare che cosa il modulo abbia deciso e perché.
\end{attenzione}

\subsection{L'identità di chi pubblica}

Il marketplace ammette due identità diverse per lo stesso artefatto: il conto
dell'utente e la macchina. «Già caricata» ha senso solo rispetto a una delle
due, ed è la ragione per cui il pannello di amministrazione conta gli
artefatti \emph{per conto e per macchina}
(\rotta{GET}{/api/admin/artefatti}), non solo in totale.
